\section{Code, data, and reproducibility}\label{appF}
\subsection{Access to code and materials}
Code and materials are provided in
\href{https://github.com/Muacca/DPPUv2-paper01}{Muacca/DPPUv2-paper01}.
The repository carries DOI
\href{https://doi.org/10.5281/zenodo.18213677}{10.5281/zenodo.18213677}.
Citations should identify the version and files used.

The symbolic core is \path{script/DPPUv2_engine_core_v3.py}.
Geometry-specific entry points are
\path{script/DPPUv2_runner_S3S1_v3.py},
\path{script/DPPUv2_runner_T3S1_v3.py}, and
\path{script/DPPUv2_runner_Nil3S1_v3.py}.
The script \path{script/DPPUv2_parameter_scan_v3.py} evaluates the
reduced potentials.
The \href{https://github.com/Muacca/DPPUv2-paper01/blob/main/script/README_ja.md}{script guide}
and \href{https://github.com/Muacca/DPPUv2-paper01/blob/main/script/docs/CONVENTIONS_ja.md}{component and global conventions}
specify usage and definitions.
Dependencies are listed in \path{script/requirements.txt}.

\subsection{Environment and calculation examples}
The following commands assume that the current directory is the
repository root. To create a virtual environment, run
\path{python -m venv .venv} and activate it in the shell being used.
Install dependencies with
\path{python -m pip install -r script/requirements.txt}.

An example symbolic calculation is:
\begin{lstlisting}[language={}]
python -B script/DPPUv2_runner_S3S1_v3.py --mode MX --ny-variant FULL --log-file new_s3_full.log
\end{lstlisting}
\Needspace{6\baselineskip}
A small scan is:
\begin{lstlisting}[language={}]
python -B script/DPPUv2_parameter_scan_v3.py --topologies S3 --ny-variants FULL --V-min 1 --V-max 4 --V-points 2 --eta-min -3 --eta-max 3 --eta-points 7 --theta-min 0 --theta-max 1 --theta-points 2 --workers 1 --output-dir new_scan
\end{lstlisting}
These inputs have $V_\tau>0$, corresponding to main-text $V<0$.
Use separate log and output destinations for each run, and retain
the association between inputs and results.

\subsection{Recording calculation conditions}
Record the code version, action variant, $T/\tau$ input convention,
volume convention, variation path, point counts and endpoints,
and classifier boundary treatment.
Distinguish $r$ from the actual sphere radius $a$, and $R_i$ from
the torus circumferences $\ell_i$.
Specify the Nil geometry including its lattice face identifications.

Appendix~\ref{appC} defines the eight CSV columns and input domains;
Appendix~\ref{appD} defines displayed quantities.
Notebook outputs store values from the execution of their respective cells.
To reproduce a calculation, specify the environment, restart the kernel,
and execute the required calculation cells in order.

\subsection{Recording and citing results}
Attach the geometry, action variant, parameter values, and code version
to results. Verify that CSV values, potential curves, and phase labels
refer to the same inputs.
Distinguish analytic results from numerical results on an input grid.
The repository's
\href{https://github.com/Muacca/DPPUv2-paper01/blob/main/LICENSE}{LICENSE}
applies. Providing the procedure, inputs, and outputs together enables
reproduction and comparison under the same conditions.
